\section*{Capítulo \ref{ch:g11:polarity-and-cohesion} --- Eletronegatividade, polaridade e forças intermoleculares}

\begin{solution}{exo:g11:polarity-and-cohesion:1}
\ce{H-Cl}: o cloro (3.16 contra 2.20). \ce{O-H}: o oxigênio (3.44).
\ce{C-O}: o oxigênio (3.44 contra 2.55).
\end{solution}

\begin{solution}{exo:g11:polarity-and-cohesion:2}
\ce{H-H}: 0, não polar. \ce{C-H}: 0.35, tratada como apolar. \ce{O-H}:
1.24, polar. \ce{N-H}: 0.84, polar. \ce{Cl-Cl}: 0, não polar. \ce{C-Cl}:
0.61, polar.
\end{solution}

\begin{solution}{exo:g11:polarity-and-cohesion:3}
Nitrogênio (3.04 > 2.19) e oxigênio (3.44 > 2.58): a \omterm{def:g11:polarity-and-cohesion:electronegativity}{eletronegatividade}
diminui descendo um grupo. Nitrogênio (3.04 > 2.55) e magnésio
(1.31 > 0.93): ela aumenta da esquerda para a direita ao longo de um
período.
\end{solution}

\begin{solution}{exo:g11:polarity-and-cohesion:4}
A amônia e a água: cada uma tem \omterm{def:g7:atoms-and-molecules:atom}{átomos} de hidrogênio ligados ao
nitrogênio ou ao oxigênio, e \omterm{def:g10:lewis-and-shape:lone-pair}{pares não ligantes} nesse \omterm{def:g7:atoms-and-molecules:atom}{átomo}. O metano
não tem hidrogênio ligado a N, O ou F; o cloreto de hidrogênio tem um
hidrogênio ligado ao cloro, que não é um dos três \omterm{def:g7:atoms-and-molecules:atom}{átomos} da definição.
\end{solution}

\begin{solution}{exo:g11:polarity-and-cohesion:5}
Sim: as \omterm{def:g11:polarity-and-cohesion:van-der-waals}{interações de van der Waals} existem entre todas as \omterm{def:g7:atoms-and-molecules:molecule}{moléculas}.
Mas as \omterm{def:g7:atoms-and-molecules:molecule}{moléculas} de metano são pequenas (10 \omterm{def:g9:inside-the-atom:nucleus}{elétrons}), por isso essas
interações são muito fracas e se rompem a uma temperatura muito baixa:
o metano ferve perto de \qty{-162}{\celsius}.
\end{solution}

\begin{solution}{exo:g11:polarity-and-cohesion:6}
\ce{CH2Cl2}: polar (duas \omterm{def:g11:polarity-and-cohesion:polar-bond}{ligações polares} \ce{C-Cl} e duas ligações
\ce{C-H} quase apolares: nenhuma simetria as anula). \ce{CCl4}: apolar
(quatro ligações idênticas nos vértices de um tetraedro). \ce{NH3}:
polar (pirâmide). \ce{CO2}: apolar (linear, ligações opostas). \ce{BF3}:
apolar (três \omterm{def:g11:polarity-and-cohesion:polar-bond}{ligações polares} idênticas a \ang{120} num plano se
anulam).
\end{solution}

\begin{solution}{exo:g11:polarity-and-cohesion:7}
O hidrogênio do grupo \ce{O-H} pode ser doado (não os do \ce{CH3},
ligados ao carbono); o \omterm{def:g7:atoms-and-molecules:atom}{átomo} de oxigênio, com os seus dois \omterm{def:g10:lewis-and-shape:lone-pair}{pares não
ligantes}, pode receber. Sim: o \ce{O-H} do metanol pode se ligar ao
oxigênio de uma \omterm{def:g7:atoms-and-molecules:molecule}{molécula} de água, e um \ce{O-H} da água ao oxigênio do
metanol, por exemplo \ce{CH3-O-H}\,$\cdots$\,\ce{OH2}.
\end{solution}

\begin{solution}{exo:g11:polarity-and-cohesion:8}
As três \omterm{def:g7:atoms-and-molecules:molecule}{moléculas} são apolares e têm a mesma geometria, mas são cada vez
maiores (10, 18 e 36 \omterm{def:g9:inside-the-atom:nucleus}{elétrons}): as suas \omterm{def:g11:polarity-and-cohesion:van-der-waals}{interações de van der Waals}
crescem, e as suas temperaturas de ebulição também.
\end{solution}

\begin{solution}{exo:g11:polarity-and-cohesion:9}
O \omterm{def:g9:periodic-table-first-look:period}{grupo} 14: o metano não tem hidrogênio ligado a N, O ou F e não forma
\omterm{def:g11:polarity-and-cohesion:hydrogen-bond}{ligações de hidrogênio}, por isso segue a tendência do seu grupo.
\end{solution}

\begin{solution}{exo:g11:polarity-and-cohesion:10}
As \omterm{def:g11:polarity-and-cohesion:van-der-waals}{interações de van der Waals}: as \omterm{def:g7:atoms-and-molecules:molecule}{moléculas} crescem do \ce{HCl} para o
\ce{HI}, e o aumento dessas interações com o tamanho supera a diminuição
da atração entre as \omterm{def:g11:polarity-and-cohesion:polar-bond}{cargas parciais}.
\end{solution}

\begin{solution}{exo:g11:polarity-and-cohesion:11}
Uma \omterm{def:g7:atoms-and-molecules:molecule}{molécula} de iodo é muito maior que uma \omterm{def:g7:atoms-and-molecules:molecule}{molécula} de cloro (106
\omterm{def:g9:inside-the-atom:nucleus}{elétrons} contra 34): as suas \omterm{def:g11:polarity-and-cohesion:van-der-waals}{interações de van der Waals} são muito mais
fortes, fortes o bastante para manter as \omterm{def:g7:atoms-and-molecules:molecule}{moléculas} num sólido à
temperatura ambiente.
\end{solution}

\begin{solution}{exo:g11:polarity-and-cohesion:12}
O propano é apolar: só \omterm{def:g11:polarity-and-cohesion:van-der-waals}{interações de van der Waals}. O metoximetano é
polar (duas \omterm{def:g11:polarity-and-cohesion:polar-bond}{ligações polares} \ce{C-O} num \ce{C-O-C} angular), mas não
tem hidrogênio no seu oxigênio: não há \omterm{def:g11:polarity-and-cohesion:hydrogen-bond}{ligações de hidrogênio} entre as
suas \omterm{def:g7:atoms-and-molecules:molecule}{moléculas}. O etanol é polar e os seus grupos \ce{O-H} formam
\omterm{def:g11:polarity-and-cohesion:hydrogen-bond}{ligações de hidrogênio}. Quanto mais fortes as atrações, mais alta a
temperatura de ebulição: propano < metoximetano < etanol.
\end{solution}

\begin{solution}{exo:g11:polarity-and-cohesion:13}
A rede aberta do gelo ocupa mais espaço que as mesmas \omterm{def:g7:atoms-and-molecules:molecule}{moléculas} no
líquido, por isso o gelo é menos denso que a água líquida e flutua. Um
lago congela de cima para baixo: a camada de gelo isola a água de baixo,
que continua líquida, e os peixes sobrevivem.
\end{solution}

\begin{solution}{exo:g11:polarity-and-cohesion:14}
Metano sólido (\omterm{def:g11:polarity-and-cohesion:van-der-waals}{interações de van der Waals} entre \omterm{def:g7:atoms-and-molecules:molecule}{moléculas} pequenas) <
gelo (\omterm{def:g11:polarity-and-cohesion:hydrogen-bond}{ligações de hidrogênio}) < cloreto de potássio (atração entre \omterm{def:g9:ions:ion}{íons}).
\end{solution}

\begin{solution}{exo:g11:polarity-and-cohesion:15}
Uma \omterm{def:g7:atoms-and-molecules:molecule}{molécula} de água tem dois \omterm{def:g7:atoms-and-molecules:atom}{átomos} de hidrogênio para doar e dois
\omterm{def:g10:lewis-and-shape:lone-pair}{pares não ligantes} com que receber: cada \omterm{def:g7:atoms-and-molecules:molecule}{molécula} pode participar de
quatro \omterm{def:g11:polarity-and-cohesion:hydrogen-bond}{ligações de hidrogênio}, duas doadas e duas recebidas, o que
constrói uma rede completa. Uma \omterm{def:g7:atoms-and-molecules:molecule}{molécula} de fluoreto de hidrogênio tem
três \omterm{def:g10:lewis-and-shape:lone-pair}{pares não ligantes}, mas só um \omterm{def:g7:atoms-and-molecules:atom}{átomo} de hidrogênio: ela doa uma
única \omterm{def:g11:polarity-and-cohesion:hydrogen-bond}{ligação de hidrogênio}, por isso, em média, só uma ligação por
\omterm{def:g7:atoms-and-molecules:molecule}{molécula} pode se formar (cadeias em zigue-zague). A água tem cerca do
dobro de \omterm{def:g11:polarity-and-cohesion:hydrogen-bond}{ligações de hidrogênio} por \omterm{def:g7:atoms-and-molecules:molecule}{molécula}, e ferve mais alto.
\end{solution}

\begin{solution}{pb:g11:polarity-and-cohesion:1}
\textbf{1.} \ce{H-O-H}, \ce{H-S-H}, \ce{H-Se-H}, cada \omterm{def:g7:atoms-and-molecules:atom}{átomo} central com
dois \omterm{def:g10:lewis-and-shape:lone-pair}{pares não ligantes}: O, S e Se têm seis \omterm{def:g10:electron-shells:valence}{elétrons de valência}, por
estarem no mesmo grupo.

\textbf{2.} $M(\ce{H2O}) = \qty{18.0}{g/mol}$, $M(\ce{H2S}) =
\qty{34.1}{g/mol}$, $M(\ce{H2Se}) = \qty{81.0}{g/mol}$.

\textbf{3.} $3.44 - 2.20 = 1.24$; $2.58 - 2.20 = 0.38$;
$2.55 - 2.20 = 0.35$. Só a ligação \ce{O-H} é claramente polar.

\textbf{4.} Angular (quatro grupos, dois deles \omterm{def:g10:lewis-and-shape:lone-pair}{pares não ligantes}). A
água é polar; o sulfeto de hidrogênio e o seleneto de hidrogênio, só
muito pouco.

\textbf{5.} 10, 18 e 36 \omterm{def:g9:inside-the-atom:nucleus}{elétrons}: o seleneto de hidrogênio tem as
\omterm{def:g11:polarity-and-cohesion:van-der-waals}{interações de van der Waals} mais fortes.

\textbf{6.} Só a água: os seus \omterm{def:g7:atoms-and-molecules:atom}{átomos} de hidrogênio estão ligados ao
oxigênio, um dos três \omterm{def:g7:atoms-and-molecules:atom}{átomos} muito eletronegativos; o enxofre e o
selênio não são eletronegativos o bastante para dar aos seus \omterm{def:g7:atoms-and-molecules:atom}{átomos} de
hidrogênio um $\delta^+$ acentuado.

\textbf{7.} Quatro: duas pelos seus próprios \omterm{def:g7:atoms-and-molecules:atom}{átomos} de hidrogênio, duas
recebidas pelos seus dois \omterm{def:g10:lewis-and-shape:lone-pair}{pares não ligantes}.

\textbf{8.} \omterm{def:g11:polarity-and-cohesion:van-der-waals}{Interações de van der Waals} (incluindo uma atração fraca
entre as pequenas \omterm{def:g11:polarity-and-cohesion:polar-bond}{cargas parciais}).

\textbf{9.} $212.9 - 273.15 = \qty{-60.3}{\celsius}$ e
\qty{-41.3}{\celsius}.

\textbf{10.} $-41.3 - (-60.3) = \qty{19.0}{\celsius}$.

\textbf{11.} A \omterm{def:g7:atoms-and-molecules:molecule}{molécula} é maior, por isso as suas \omterm{def:g11:polarity-and-cohesion:van-der-waals}{interações de van der
Waals} são mais fortes.

\textbf{12.} $-60.3 - 19.0 = \qty{-79.3}{\celsius}$.

\textbf{13.} Subida $-88.1 - (-112) = \qty{23.9}{\celsius}$; previsão
para o metano $-112 - 23.9 = \qty{-136}{\celsius}$, contra
\qty{-162}{\celsius} na realidade: a extrapolação fica cerca de
\qty{26}{\celsius} alta demais. Ela não é exata; a tendência real não é
uma reta.

\textbf{14.} Subida $\qty{18.7}{\celsius}$; previsão para o fluoreto de
hidrogênio $-85.1 - 18.7 = \qty{-103.8}{\celsius}$, contra
\qty{19.6}{\celsius}: uma diferença de cerca de \qty{123}{\celsius}.

\textbf{15.} $100 - (-79.3) = \qty{179}{\celsius}$, cerca de
\qty{180}{\celsius}.

\textbf{16.} Subida $\qty{25.3}{\celsius}$; previsão $-87.8 - 25.3 =
\qty{-113}{\celsius}$; diferença $-33.4 - (-113) = \qty{80}{\celsius}$.

\textbf{17.} Água (\qty{180}{\celsius}) > fluoreto de hidrogênio
(\qty{123}{\celsius}) > amônia (\qty{80}{\celsius}). A água doa dois
\omterm{def:g7:atoms-and-molecules:atom}{átomos} de hidrogênio e recebe com dois \omterm{def:g10:lewis-and-shape:lone-pair}{pares não ligantes}: quatro
\omterm{def:g11:polarity-and-cohesion:hydrogen-bond}{ligações de hidrogênio} por \omterm{def:g7:atoms-and-molecules:molecule}{molécula}, todas usadas. O fluoreto de
hidrogênio só tem um hidrogênio para doar, e a amônia só um \omterm{def:g10:lewis-and-shape:lone-pair}{par não
ligante} com que receber: em média, menos ligações por \omterm{def:g7:atoms-and-molecules:molecule}{molécula}, e o
nitrogênio, menos eletronegativo, forma ligações mais fracas.

\textbf{18.} Um gás: ela ferveria perto de \qty{-80}{\celsius}. A Terra
não teria água líquida, nem oceanos, nem chuva, nem nenhuma das formas
de vida que dependem deles.

\textbf{19.} Ela extrapola uma reta traçada por apenas dois pontos, e o
teste no \omterm{def:g9:periodic-table-first-look:period}{grupo} 14 mostra que uma extrapolação assim pode errar em uns
\qty{25}{\celsius}.

\textbf{20.} Sem \omterm{def:g11:polarity-and-cohesion:hydrogen-bond}{ligações de hidrogênio}, a água ferveria a cerca de
\qty{-80}{\celsius}, uns \qty{180}{\celsius} abaixo do seu verdadeiro
ponto de ebulição.
\end{solution}
